% Mark Gilbert: Narrating the Process – Reading Summary
% Introduction to the European Union, Week 5
% Compile with: pdflatex gilbert-narrating-the-process-reading-summary.tex

\documentclass[11pt,a4paper]{article}
\usepackage[utf8]{inputenc}
\usepackage[T1]{fontenc}
\usepackage[margin=2.5cm]{geometry}
\usepackage{booktabs}
\usepackage{enumitem}
\setlist{nosep,leftmargin=*}
\usepackage{parskip}
\usepackage{microtype}
\usepackage{hyperref}
\usepackage{xcolor}
\definecolor{eublue}{RGB}{0,51,153}

\hyphenation{Spitzenkandidat Spitzenkandidaten}
\hypersetup{
  colorlinks=true,
  linkcolor=eublue,
  urlcolor=eublue,
  pdftitle={Gilbert – Narrating the Process},
  pdfauthor={Introduction to the EU, UCM}
}

\begin{document}

\title{Mark Gilbert: Narrating the Process\\Questioning the Progressive Story of European Integration}
\author{Introduction to the European Union\\Universidad Complutense de Madrid\\Week 5\\\textit{Journal of Common Market Studies} (2008)}
\date{}
\maketitle
\vspace{-1em}
\hrule
\vspace{1em}

\section*{Rhetorical Framework (Ethos, Logos, Pathos)}

\begin{table}[h]
\centering
\begin{tabular}{@{}lll@{}}
\toprule
\textbf{Appeal} & \textbf{Meaning} & \textbf{Gilbert's Use} \\
\midrule
\textbf{Ethos} & Credibility, authority, trust & Historiographical critic; challenges dominant narratives; advocates for methodological reform \\
\textbf{Logos} & Logic, structure, argument & Critique of teleology; contingency, dissent, losers as analytical categories \\
\textbf{Pathos} & Emotion, stakes, persuasion & Appeal to readers who distrust ``success story'' histories; urgency of honest engagement \\
\bottomrule
\end{tabular}
\end{table}

\textbf{Key insight:} Gilbert uses \textit{logos} (teleology as problem, contingency as alternative) to build \textit{ethos} (as critical historian). His \textit{pathos} targets readers who feel the official EU story is too smooth---crisis, power, exclusion deserve more space.

\section*{Bibliographic Information}

\begin{itemize}
\item \textbf{Author:} Mark Gilbert
\item \textbf{Title:} Narrating the Process: Questioning the Progressive Story of European Integration
\item \textbf{Source:} Journal of Common Market Studies (JCMS), Vol. 46, No. 3
\item \textbf{Year:} 2008
\item \textbf{Pages:} 641--662
\end{itemize}

\section*{Summary \& Key Points}

\begin{itemize}
\item Gilbert criticises the dominant ``progressive'' narrative that portrays European integration as a linear story of peace, prosperity and democracy.
\item He argues that such teleological accounts downplay conflict, failure and alternative paths that were possible at each moment.
\item The article calls for more critical histories that treat the EU as one contested political project among many, rather than as the inevitable end-point of European history.
\end{itemize}

\section*{Main Arguments}

\subsection*{I. Normative, Teleological Histories}

Much writing about European integration is implicitly normative: it is a history \textit{for} Europe, justifying the EU by projecting a coherent, upward trajectory from war to peace. Gilbert argues that this obscures the messy, contingent nature of the process.

\subsection*{II. Foregrounding Contingency, Dissent, and Losers}

He proposes that historians should foreground contingency, dissent and losers of integration, giving more space to actors and ideas that do not fit the official success story. This would make our understanding of the EU more realistic and less celebratory.

\section*{Context \& Analysis}

``Questioning'' in the title = \textit{logos} (interrogation as argument) + \textit{ethos} (Gilbert as critic). ``Progressive story'' = target of critique; \textit{pathos} in what gets suppressed (conflict, failure, alternatives).

\section*{Relevance to Course}

\begin{itemize}
\item \textbf{Ethos:} Meta-level framework for evaluating how EU history is told in textbooks, speeches, official narratives.
\item \textbf{Logos:} Avoid teleology in your own essays---discuss crises, failures, alternatives without assuming inevitability.
\item \textbf{Pathos:} Connects to legitimacy debates---if the success story is questioned, new justifications (or criticisms) are needed.
\end{itemize}

\section*{Discussion Questions}

\begin{itemize}
\item Dangers of presenting integration as inevitable success story? (Ethos: who benefits? Logos: what gets obscured? Pathos: stakes of legitimacy)
\item How might historians integrate colonialism, dissent, conflict? (Ethos: whose voices? Logos: methodological choices; Pathos: inclusion, exclusion)
\item Do political actors need simplified stories even if historians know they're incomplete? (Ethos: politicians vs.\ historians; Logos: trade-off: motivation vs.\ accuracy; Pathos: legitimacy, public support)
\end{itemize}

\vspace{1em}
\noindent\footnotesize\textit{Reference:} Gilbert, Mark (2008). ``Narrating the Process: Questioning the Progressive Story of European Integration.'' \textit{Journal of Common Market Studies}, 46(3), pp. 641--662.

\end{document}
